% GENERATED FILE. Edit canonical lessons, then run npm run generate:books.
% canonical-source-hash: fnv1a64:c7dbb2bd7dd52785
% canonical-lessons: ES-C45-repaso-ocho

\chapter{Review --- All Eight}
\label{ch:review-object-complete}
\hlchaptermodality{91}

\begin{chapteropening}
\textbf{By the end of this chapter:} \emph{I can give all eight direct object pronouns and say which four depend on a noun.}

Complete ES-C45-repaso-ocho: give the eight pronouns and identify the four that a noun decides.
\end{chapteropening}

\section[Practice]{Review — all eight}

\label{lesson:ES-C45-repaso-ocho}



\begin{quote}
Nothing new here. Before the grid: where does an object pronoun sit, and has that ever changed?
\end{quote}

\begin{grammarlens}[title={What you've built}]
Never, in nine chapters. Here is the complete set:

\noindent
\begin{tabularx}{\linewidth}{@{}>{\raggedright\arraybackslash}X>{\raggedright\arraybackslash}X>{\raggedright\arraybackslash}X@{}}
\toprule
\textbf{} & \textbf{singular} & \textbf{plural} \\
\midrule
\textbf{me} & \textbf{me} & \textbf{nos} \\
\textbf{you} & \textbf{te} & \textbf{os} \emph{(Spain)} \\
\textbf{it / them} — masculine & \textbf{lo} & \textbf{los} \\
\textbf{it / them} — feminine & \textbf{la} & \textbf{las} \\
\bottomrule
\end{tabularx}

Eight cells, built one at a time, and this is the first page on which they have all appeared together. That is deliberate. A grid handed to you on the day you met the first one would have been eight things to memorise; here it is a picture of ground you already hold.

Count the genuinely new \textbf{shapes} you had to learn across those chapters: \emph{me} was already yours from chapter four, \emph{te} was one, \emph{nos} one, \emph{os} one, \emph{lo} one, \emph{la} one — and \emph{los} and \emph{las} were free, because you produced them yourself from an \emph{-s}.
\end{grammarlens}

\begin{grammarlens}[title={Grammar Lens: two questions}]
Reading the grid as eight decisions is the hard way. It is two:

1. \textbf{Who is on the receiving end} — me, you, or something else? 2. \textbf{How many, and if it is a thing, which gender?}

The second question only bites in the bottom two rows, because only there does a noun decide. \emph{Me}, \emph{te}, \emph{nos} and \emph{os} are settled by the conversation itself: you always know who is talking.

So the difficulty is not eight-eighths of this table. It is one quarter of it.
\end{grammarlens}

\subsection*{Your turn}
Say these aloud:

\begin{itemize}
\raggedright
  \item the eight in grid order, twice
  \item \textquotedblleft{}me quieres, te quiero, nos quieren, os quiero\textquotedblright{}
\end{itemize}

\emph{Twice through:} Until the grid is a reflex rather than a lookup.

\begin{tcolorbox}[breakable,colback=teal!4,colframe=teal!35!black,title={Before you move on}]
How many direct object pronouns are there? (\textbf{Eight}.) How many of them depend on a noun? (\textbf{Four}.) Next: what all eight are for, in a conversation that would be unbearable without them.
\end{tcolorbox}
